% SDKP_manuscript.tex
\documentclass[10pt,twocolumn]{article}
\usepackage[utf8]{inputenc}
\usepackage{amsmath,amssymb}
\usepackage{graphicx}
\usepackage{caption}
\usepackage{authblk}
\usepackage{hyperref}
\usepackage{cite}

\title{The SDKP Framework: Time–Density–Kinematics as a Unified Mechanism for Gravity and Quantum Phenomena}
\author{Donald Paul Smith (FatherTimeSDKP)\\
\small Independent Researcher}
\date{June 7, 2026}

\begin{document}
\twocolumn[
\maketitle
\begin{abstract}
SDKP proposes a unified ontology where Scale (S), Density (D), and Kinematic flux (K) are primary field variables. Gravity and inertial effects emerge from spatial density gradients and kinematic coupling; time corresponds to information/entropy flux of these variables; quantum behavior arises from lattice information dynamics. We present core postulates, principal equations, limiting cases that recover Newtonian and quantum-like behavior, and a set of falsifiable predictions with proposed validation paths.
\end{abstract}
]

\section*{Introduction}
The SDKP framework posits three primary fields: scale $S(x,t)$, density $D(x,t)$, and kinematic flux $K(x,t)$; these determine emergent observables conventionally associated with mass, time, and quantum phase. The motivation is to supply a minimal ontology that produces gravitational and quantum phenomenology from the same underlying dynamics.

\section*{Postulates}
1. Existence of continuous (or fine-grained discrete) SDKP fields $S,D,K$.\\
2. Local conservation: $\partial_t D + \nabla\cdot K = 0$.\\
3. Scale–density coupling: local pressure-like functional $P(D,S)$ governs emergent forces.\\
4. Time-rate as information flux $\tau^{-1}\propto\Phi(D,K,\nabla S)$ (schematic).\\
5. Quantum amplitudes arise from phase-encoded lattice operator algebra on cells of characteristic scale $\ell_S$.

\section*{Core equations (schematic)}
Conservation:
\begin{equation}
\partial_t D + \nabla\cdot K = 0.
\end{equation}
Kinematic evolution (phenomenological):
\begin{equation}
\partial_t K + \nabla\cdot (K\otimes v) = -\nabla P(D,S) + \mathcal{C}(D,K,S),
\end{equation}
with $v$ a local velocity-like field and $\mathcal{C}$ coupling terms (rotation, dissipation). Emergent force on test element:
\begin{equation}
F_{\rm emergent} \approx -\nabla P(D,S).
\end{equation}
Time-rate relation (schematic):
\begin{equation}
\frac{d\tau^{-1}}{dt} = \alpha\,\Phi(D,K,\omega),
\end{equation}
where $\omega$ denotes local rotation and $\alpha$ a scaling constant.

\section*{Recovery of known limits}
Newtonian limit: for slowly varying $S,D$ and weak gradients, $P(D,S)\sim G_{\rm eff} M_{\rm eff} / r^2$ yields inverse-square behavior with effective gravitational constant $G_{\rm eff}$ derived from microparameters. Relativistic corrections: density+rotation terms modify proper-time accumulation, producing corrections to standard time-dilation. Quantum limit: lattice operator algebra produces linear superposition-like behavior for cell amplitudes; coarse-graining gives Schr\"odinger-like evolution with decoherence controlled by coupling to macroscopic density modes.

\section*{Representative derivation (concise)}
Consider single-axis radial symmetry and small perturbations $\delta D$ about background $D_0$. Linearizing the conservation and kinematic equations and applying ansatz $K\sim D v$ yields wave-like modes with dispersion relation:
\begin{equation}
\omega^2 = c_S^2 k^2 + \beta k^4 + \dots
\end{equation}
Low-$k$ behavior reproduces massless propagation; higher-order terms introduce localization and bound-state-like modes analogous to quantum stationary states when projected onto lattice eigenmodes.

\section*{Predictions and observables}
- Schumann resonance centroid shifts: global density–rotation events produce measurable Δf in ELF modal frequencies (claimed order $10^{-3}$–$10^{-2}$ Hz).\\
- Precision clock anomalies: density+rotation correction to time-dilation near strong density gradients.\\
- Mesoscopic decoherence scaling: predicted dependence of coherence length on local density modes in conductive samples.\\
- Laboratory superconducting collapse under controlled density–rotation coupling parameters.

\section*{Proposed validation path}
1. Reproduce lattice simulations and analytical derivations.\\
2. Analyze existing global ELF datasets for predicted Schumann signatures.\\
3. Design targeted laboratory experiments for superconducting collapse and mesoscopic decoherence tests.\\
4. Publish incremental, narrowly scoped results for community scrutiny.

\section*{Reproducibility}
Simulation code and datasets should be released under open license; minimal example reproducing the lattice bound-state behavior is essential.

\section*{Discussion}
SDKP is an alternative framework with explicit, falsifiable predictions. Primary avenues for acceptance are independent reproduction, successful prediction of an unexplained phenomenon, or consistent constraints that falsify core claims.

\section*{Conclusion}
We have outlined the SDKP postulates, core equations, limiting behavior, and experimental pathways. Focused replication and targeted ELF and laboratory tests are recommended to evaluate the framework's viability.

\section*{Acknowledgements}
Author acknowledges feedback from independent reviewers and lists repositories for code and data (placeholders).

\bibliographystyle{unsrt}
\begin{thebibliography}{9}
\bibitem{sdkp1} D. P. Smith, ``SDKP Manuscript,'' self-published archive (placeholder), 2025.
\bibitem{sdkp2} D. P. Smith, ``Simulations and datasets,'' GitHub/Zenodo (placeholder), 2025.
\end{thebibliography}

\end{document}
